\section{第IVポジション}
\begin{center}
\begin{tabular}{|lcl|}
\hline
この章の基礎練習 & : & 開放弦の練習\\ 
この章の修了課題 & : & 1. 「\ref{half_scale}」の音階練習を正しい音程で暗譜して演奏できる\\
               &   & 2. ベートーヴェン、チャイコフスキー、「山の魔王の宮殿にて」からの抜粋を正しい音程で暗譜して演奏できる\\
\hline
\end{tabular}
\end{center}
\subsection{ベーシストの左手をつくる}
\begin{flushleft}
\begin{minipage}{203pt}
\ \ \ \ 左手の指には図\addtocounter{figure}{1}\thefigure のように番号が付いています。音符の上に書いてある数字に対応する指で弦を押さえます。例えば、4と書いてあったら小指で弦を押さえます。開放弦は0で表します。3の指は第\ref{6th}章ま
で使いません。\\
\end{minipage}
\hfill
\begin{minipage}{110pt}
\addtocounter{figure}{-1}
\begin{center}
\includegraphics[height=3.6cm]{../Vol1/Pics/Position/fingering.epsi}\\
{\small 図\thefigure : 指番号と縛る指\\}
\end{center}
\end{minipage}
\hfill
\begin{minipage}{120pt}
\addtocounter{figure}{1}
\begin{center}
\includegraphics[height=3.6cm]{../Vol1/Pics/photo0830/open_finger.epsi}\\
{\small 図\thefigure : ネックの根元に差し込む\\}
\end{center}
\end{minipage}
\end{flushleft}

左手の練習に入る前に、左手の2と3の指が離れないように縛ってください(図
\thefigure)。そして、1と2の間、2と4の間の距離が等しくなるように指を開
きます。左手の形が身に付くまでは図\thefigure のように指を縛ってから練
習します。比較的開きにくい1と2の指の間を広げるには、図
\addtocounter{figure}{1}\thefigure のようにネックの根元に指の間を差し
込むことを続けると有効です。\\

\subsection{第IVポジションの位置}
\begin{flushleft}
\begin{minipage}{160pt}
　最初に紹介する第IVポジションはネックの根元にあるポジションです(図
\addtocounter{figure}{2}\thefigure 、 
\addtocounter{figure}{1}\thefigure)。親指をネックの一番下まで降ろして
ください。その親指の正面を2の指で押さえます。これが第IVポジションです。
1と4は2に合わせます(図\addtocounter{figure}{1}\thefigure)。
\end{minipage}
\hfill
\begin{minipage}{80pt}
\addtocounter{figure}{-2}
\begin{center}
%\includegraphics[height=4.3cm]{../Vol1/Pics/Position/half_1.epsi}\\
\includegraphics[width=2cm]{../Vol1/Pics/now_printing.eps}\\
\end{center}
{\small 図\thefigure : 第IVポジションの位置\\}
\end{minipage}
\hfill
\begin{minipage}{80pt}
\addtocounter{figure}{1}
\begin{center}
\includegraphics[height=6.5cm]{../Vol1/Pics/Position/4th_2.epsi}\\
\end{center}
{\small 図\thefigure : 奏者から見た第IVポジション}\\
\end{minipage}
\hfill
\begin{minipage}{85pt}
\addtocounter{figure}{1}
\begin{center}
\includegraphics[height=4.3cm]{../Vol1/Pics/photo0830/left_thumb1.epsi}\\
\end{center}
{\small 図\thefigure : 左手親指の位置}\\
\end{minipage}
\end{flushleft}

\subsection{弦の押さえ方}
\indent 弦を押さえる際には1と2、2と4の指の間が均等に開くようにし、指の
先で押さえます(図\addtocounter{figure}{-2}\thefigure)。1の指はネックの
上側に少し傾け、指のやや側面寄りの位置で押さえます(図
\addtocounter{figure}{3}\thefigure)。 指は丸めてアーチ状にします(図
\thefigure)。図\addtocounter{figure}{1}\thefigure のようにつぶれたアー
チは正しくありません。
\begin{flushleft}
\begin{minipage}{220pt}
　また、2、4の指で押さえる場合には必ず上側の指も一緒に押さえます
(図\addtocounter{figure}{1}\thefigure)。たとえば2の指を押さえる際に
1を放してはいけません(図\addtocounter{figure}{1}\thefigure)。\\

\begin{minipage}{100pt}
\addtocounter{figure}{-3}
\begin{center}
\includegraphics[height=2cm]{../Vol1/Pics/newphoto/arch1.epsi}\\
\end{center}
{\small 図\thefigure : アーチ状に丸める。1は指の側面で押さえる。\\}
\end{minipage}
\hfill
\begin{minipage}{100pt}
\addtocounter{figure}{1}
\begin{center}
\includegraphics[height=2cm]{../Vol1/Pics/newphoto/arch2.epsi}\\
\end{center}
{\small 図\thefigure : アーチがつぶれてはダメ\\}
\end{minipage}
\end{minipage}
\hfill
\begin{minipage}{100pt}
\addtocounter{figure}{1}
\begin{center}
\includegraphics[width=3cm]{../Vol1/Pics/newphoto/finger1.epsi}\\
\end{center}
{\small 図\thefigure : 1の指も一緒に\\}
\end{minipage}
\hfill
\begin{minipage}{100pt}
\addtocounter{figure}{1}
\begin{center}
\includegraphics[width=3cm]{../Vol1/Pics/newphoto/finger2.epsi}\\
\end{center}
{\small 図\thefigure : 1を放してはダメ\\}
\end{minipage}
\end{flushleft}

\subsection{肘の位置}
\begin{flushleft}
\begin{minipage}{250pt}
　弦の押さえ方を覚えたら、次は肘の位置を習得しましょう。前腕(手首から肘まで)の角度が、ネックの上下方向と直角になる位置まで肘を上げます(図\addtocounter{figure}{1}\thefigure)。\\

\ \ \ \ 4の指は他の指と比べて押さえる力が弱いですが、このように肘を上げることで4の指の力を腕の重みで補えるようになります。
\end{minipage}
\hfill
\begin{minipage}{180pt}
\begin{center}
\includegraphics[height=4.5cm]{../Vol1/Pics/Shifting/orthogonal_halfsize.epsi}\\
図\thefigure : 腕とネックが直角になるのが目安\\
\end{center}
\end{minipage}
\end{flushleft}

\subsection{第IVポジションで取れる音}
第IVポジションではそれぞれの弦で以下の音を取ることができます。正し
い場所を押さえているかどうか、チューニングメーターで確認しながら弾いてみましょう。\\

\begin{music}
\nostartrule
\parindent 0pt
\setclef1{\bass}  
\startpiece
\notes\enotes
\Notes\zchar{18}{G線}\zchar{13}{\bf 1}\wh{d}\zchar{14}{\bf 2}\wh{^d}\zchar{14}{\bf 4}\wh{e}\enotes
\doublebar
\Notes\zchar{13}{\bf 1}\wh{d}\zchar{14}{\bf 2}\wh{_e}\zchar{14}{\bf 4}\wh{=e}\enotes
\doublebar
\Notes\zchar{18}{D線}\zchar{11}{\bf 1}\wh{a}\zchar{11}{\bf 2}\wh{^a}\zchar{11}{\bf 4}\wh{b}\enotes
\doublebar
\Notes\zchar{11}{\bf 1}\wh{a}\zchar{11}{\bf 2}\wh{_b}\zchar{11}{\bf 4}\wh{=b}\enotes
\setdoublebar
\endpiece
\startpiece
\notes\enotes
\Notes\zchar{14}{A線}\zchar{9}{\bf 1}\wh{'E}\zchar{9}{\bf 2}\wh{F}\zchar{9}{\bf 4}\wh{^F}\enotes
\doublebar
\Notes\zchar{9}{\bf 1}\wh{'E}\zchar{9}{\bf 2}\wh{F}\zchar{9}{\bf 4}\wh{_G}\enotes
\doublebar
\Notes\zchar{14}{E線}\zchar{9}{\bf 1}\wh{'B}\zchar{9}{\bf 2}\wh{C}\zchar{9}{\bf 4}\wh{^C}\enotes
\doublebar
\Notes\zchar{9}{\bf 1}\wh{'B}\zchar{9}{\bf 2}\wh{C}\zchar{9}{\bf 4}\wh{_D}\enotes
\setdoublebar
\endpiece
\end{music}

\subsection{音階練習 \label{half_scale}}
\begin{music}
\nostartrule
\parindent 0pt
\setclef1{\bass}  

\generalmeter{\meterC}
\generalsignature{0}    
\startpiece
\notes\zchar{14}{ハ長調(C-Dur)音階}\enotes
\NOtes\zchar{9}{\bf 2}\qu{'C}\zchar{9}{\bf 0}\ql{D}\zchar{9}{\bf 1}\ql{E}\zchar{9}{\bf 2}\ql{F}\enotes
\bar
\NOtes\zchar{9}{\bf 0}\ql{'G}\zchar{10}{\bf 1}\ql{!a}\zchar{11}{\bf 4}\ql{b}\qp\enotes
\bar
\NOtes\zchar{11}{\bf 4}\ql{b}\zchar{10}{\bf 1}\ql{a}\zchar{9}{\bf 0}\ql{'G}\zchar{9}{\bf 2}\ql{F}\enotes
\bar
\NOtes\zchar{9}{\bf 1}\ql{'E}\zchar{9}{\bf 0}\ql{D}\zchar{9}{\bf 2}\qu{C}\qp\enotes
\setdoublebar\endpiece

\generalmeter{\meterC}
\generalsignature{0}    
\startpiece
\notes\zchar{14}{イ短調(a-moll)音階}\enotes
\NOtes\zchar{9}{\bf 0}\qu{'A}\zchar{9}{\bf 1}\qu{B}\zchar{9}{\bf 2}\qu{C}\zchar{9}{\bf 0}\ql{D}\enotes
\bar
\NOtes\zchar{9}{\bf 1}\ql{'E}\zchar{9}{\bf 2}\ql{F}\zchar{10}{\bf 0}\ql{G}\zchar{11}{\bf 1}\ql{!a}\enotes
\bar
\NOtes\zchar{11}{\bf 1}\ql{a}\zchar{10}{\bf 0}\ql{'G}\zchar{9}{\bf 2}\ql{F}\zchar{9}{\bf 1}\ql{E}\enotes
\bar
\NOtes\zchar{9}{\bf 0}\ql{'D}\zchar{9}{\bf 2}\qu{C}\zchar{9}{\bf 1}\qu{B}\zchar{9}{\bf 0}\qu{A}\enotes
\setdoublebar\endpiece
\end{music}

\subsection{第IVポジションだけで弾ける名曲}
\subsubsection*{ワーグナー: 歌劇「さまよえるオランダ人」序曲より}
\begin{music}
\nostartrule
\setclef1{\bass}
\generalsignature{-1}    
\generalmeter{\meterfrac64}
\parindent 0pt
\startbarno=129
\def\writebarno{\tenrm\the\barno\barnoadd}
\def\raisebarno{2\internote}
\def\shiftbarno{0.1\Interligne}
\systemnumbers
\startpiece\bigaccid
\notes\zchar{18}{\bf Molto animato}\enotes
\Notes\zchar{12}{\upbow}\zchar{9}{\bf 0}\qu{'A}\zchar{-5}{\ff}\enotes
\bar
\Notes\zchar{-5}{\it marcato}\zchar{12}{\downbow}\zchar{9}{\bf 0}\hlp{'D}\zchar{9}{\upbow}\hl{D}\zchar{9}{\upbow}\qu{A}\enotes
\bar
\Notes\zchar{12}{\downbow}\zchar{9}{\bf 0}\hl{'D}\zchar{9}{\upbow}\ql{D}\zchar{15}{\downbow}\zchar{12}{\bf 1}\isluru{0}{!a}\hlp{a}\enotes\barno=130
\bar
\Notes\tslur{0}{a}\hl{a}\zchar{12}{\upbow}\zchar{9}{\bf 0}\ql{'D}\zchar{9}{\downbow}\hl{D}\zchar{9}{\upbow}\ql{D}\enotes
\bar
\Notes\zchar{17}{\downbow}\zchar{14}{\bf 1}\isluru{0}{d}\whp{d}\enotes
\bar
\Notes\tslur{0}{d}\ql{d}\zchar{12}{\upbow}\zchar{9}{\bf 0}\ql{'D}\zchar{9}{\downbow}\ql{D}\zchar{9}{\upbow}\ql{D}\zchar{9}{\downbow}\ql{D}\zchar{9}{\upbow}\ql{D}\enotes
\bar
\Notes\zchar{12}{\downbow}\zchar{9}{\bf 0}\ql{'G}\enotes
\setdoublebar
\endpiece
\end{music}

\documentclass{jarticle}
\usepackage{musixdoc}
\startmuflex\makeindex

\begin{document}

\subsubsection*{�쥹�ԡ���: ������֥����ޤξ��פ��֥���������ն�ξ���}
\begin{music}
\nostartrule
\setclef1{\bass}
\generalsignature{1}    
\generalmeter{\meterfrac54}
\parindent 0pt
\startbarno=28
\def\writebarno{\tenrm\the\barno\barnoadd}
\def\raisebarno{2\internote}
\def\shiftbarno{0.1\Interligne}
\systemnumbers
\startpiece\bigaccid
\notes\zchar{18}{\bf Ancora pi\`{u} mosso (\metron{\qu}{69})}\enotes
\notes\ibl{0}{b}{0}\ibbl{0}{b}{0}\qb{0}{b}\qb{0}{b}\qb{0}{b}\tbbl{0}\tbl{0}\qb{0}{b}\enotes
\notes\ibl{0}{b}{-2}\ibbl{0}{b}{-2}\qb{0}{b}\qb{0}{b}\qb{0}{b}\tbbl{0}\tbl{0}\qb{0}{a}\enotes
\notes\ibl{0}{b}{0}\ibbl{0}{b}{0}\qb{0}{b}\qb{0}{b}\qb{0}{d}\tbbl{0}\tbl{0}\qb{0}{b}\enotes
\notes\ibl{0}{b}{-2}\qb{0}{b}\ibbl{0}{b}{-2}\qb{0}{a}\tbbl{0}\tbl{0}\qb{0}{a}\enotes
\Notes\ql{b}\enotes
\bar
\Notes\hlp{b}\hl{b}\enotes
\bar
\notes\ibl{0}{b}{0}\ibbl{0}{b}{0}\qb{0}{b}\qb{0}{b}\qb{0}{b}\tbbl{0}\tbl{0}\qb{0}{b}\enotes
\notes\ibl{0}{b}{-2}\ibbl{0}{b}{-2}\qb{0}{b}\qb{0}{b}\qb{0}{b}\tbbl{0}\tbl{0}\qb{0}{a}\enotes
\notes\ibl{0}{b}{0}\ibbl{0}{b}{0}\qb{0}{b}\qb{0}{b}\qb{0}{d}\tbbl{0}\tbl{0}\qb{0}{b}\enotes
\notes\ibl{0}{b}{-2}\qb{0}{b}\ibbl{0}{b}{-2}\qb{0}{a}\tbbl{0}\tbl{0}\qb{0}{a}\enotes
\Notes\ql{b}\enotes
\bar
\Notes\hlp{b}\hl{b}\enotes
\bar
\notes\ibl{0}{b}{0}\ibbl{0}{b}{0}\qb{0}{b}\qb{0}{b}\qb{0}{b}\tbbl{0}\tbl{0}\qb{0}{b}\enotes
\notes\ibl{0}{b}{-2}\ibbl{0}{b}{-2}\qb{0}{b}\qb{0}{b}\qb{0}{b}\tbbl{0}\tbl{0}\qb{0}{a}\enotes
\notes\ibl{0}{b}{0}\ibbl{0}{b}{0}\qb{0}{b}\qb{0}{b}\qb{0}{d}\tbbl{0}\tbl{0}\qb{0}{b}\enotes
\notes\ibl{0}{b}{-2}\qb{0}{b}\ibbl{0}{b}{-2}\qb{0}{a}\tbbl{0}\tbl{0}\qb{0}{a}\enotes
\Notes\ql{b}\enotes 
\bar
\Notes\hlp{b}\hl{b}\enotes
\setdoublebar
\endpiece
\end{music}
\endmuflex
\end{document}

\subsubsection*{マーラー: 交響曲第1番 ニ長調 第2楽章より}

\begin{music}
\nostartrule
\setclef1{\bass}
\generalsignature{3}
\generalmeter{\meterfrac34}
\parindent 0pt
\def\writebarno{\tenrm\the\barno\barnoadd}
\def\raisebarno{2\internote}
\def\shiftbarno{0.1\Interligne}
\startbarno=285
\systemnumbers
\startpiece\bigaccid
\notes\zchar{17}{\Bigtype 26}\sk  \zchar{19}{\ \ \bf Tempo primo. (Kr\"{a}ftig bewegt, doch nicht zu schunell)}\enotes
\notes\loffset{1.5}{\zchar{-8}{\ff}}\zchar{12}\downbow\zchar{10}{\bf 1}\hl{a}\zchar{-5}{\bf 1}\sk\enotes
\Notes\zchar{12}\upbow\zchar{9}{\bf 1}\upz{'E}\ql{E}\enotes
\bar
\Notes\zchar{14}\downbow\zchar{11}{\bf 1}\upz{a}\ql{a}\zchar{14}\upbow\zchar{11}{\bf 1}\upz{a}\ql{a}\zchar{11}\upbow\zchar{8}{\bf 1}\upz{'E}\ql{E}\enotes
\bar
\notes\hl{a}\zchar{-5}{\bf 2}\sk\enotes
\Notes\upz{'E}\ql{E}\enotes
\bar
\Notes\upz{a}\ql{a}\upz{a}\ql{a}\upz{'E}\ql{E}\enotes
\bar
\notes\hl{a}\zchar{-5}{\bf 3}\sk\enotes
\Notes\upz{'E}\ql{E}\enotes
\bar
\Notes\upz{a}\ql{a}\upz{a}\ql{a}\upz{'E}\ql{E}\enotes
\bar
\notes\hl{a}\zchar{-5}{\bf 4}\sk\enotes
\Notes\upz{'E}\ql{E}\enotes
\bar
\Notes\upz{a}\ql{a}\upz{a}\ql{a}\zchar{-5}{\p}\upz{'E}\ql{E}\enotes
\bar
\notes\hl{a}\zchar{-5}{\bf 5}\sk\enotes
\Notes\upz{'E}\ql{E}\enotes
\bar
\Notes\upz{a}\ql{a}\upz{a}\ql{a}\zchar{-5}{\p}\upz{'E}\ql{E}\enotes
\bar
\notes\zchar{9}\downbow\hl{a}\sk\enotes
\Notes\zchar{9}\upbow\ql{'E}\enotes
\bar
\Notes\zchar{10}\downbow\upz{a}\ql{a}\zchar{10}\upbow\upz{a}\ql{a}\zchar{10}\downbow\upz{a}\ql{a}\enotes
\bar
\Notes\zchar{15}\upbow\zchar{12}{\bf 4}\upz{b}\ql{b}\zchar{15}\downbow\zchar{12}{\bf 4}\upz{b}\ql{b}\zchar{13}\upbow\zchar{10}{\bf 1}\lpz{'B}\qu{B}\enotes
\bar
\Notes\ql{'E}\qp\qp\enotes
\setdoublebar
\endpiece

\nostartrule
\setclef1{\bass}
\generalsignature{3}
\generalmeter{\meterfrac34}
\parindent 0pt
\def\writebarno{\tenrm\the\barno\barnoadd}
\def\raisebarno{2\internote}
\def\shiftbarno{0.1\Interligne}
\startbarno=303
\systemnumbers
\startpiece\bigaccid
\notes\zchar{12}{\Bigtype 27}\zchar{-8}{\mf}\enotes
\notes\zchar{9}{\downbow}\zchar{-4}{\bf 0}\hu{'A}\zchar{-8}{\bf 1}\sk\enotes
\Notes\zchar{9}\upbow\zchar{-7}{\bf 0}\lpz{E}\qu{E}\enotes
\bar
\Notes\zchar{9}\downbow\lpz{'A}\qu{A}\zchar{9}\upbow\lpz{A}\qu{A}\zchar{9}\upbow\lpz{!E}\qu{!E}\enotes
\bar
\notes\hu{'A}\zchar{-6}{\bf 2}\sk\enotes
\Notes\lpz{E}\qu{E}\enotes
\bar
\Notes\lpz{'A}\qu{A}\lpz{A}\qu{A}\lpz{!E}\qu{E}\enotes
\bar
\notes\hu{'A}\zchar{-6}{\bf 3}\sk\enotes
\Notes\lpz{E}\qu{E}\enotes
\bar
\Notes\lpz{'A}\qu{A}\lpz{A}\qu{A}\lpz{!E}\qu{E}\enotes
\bar
\notes\hu{'A}\zchar{-6}{\bf 4}\sk\enotes
\Notes\lpz{E}\qu{E}\enotes
\bar
\Notes\zchar{10}{\downbow}\lpz{'A}\qu{A}\zchar{10}{\upbow}\lpz{A}\qu{A}\zchar{10}{\downbow}\lpz{!E}\qu{E}\enotes
\bar
\Notes\zchar{10}{\upbow}\lsfz{'B}\hup{B}\enotes
\bar
\Notes\loffset{2.8}{\zchar{13}{\Bigtype 28}}\zchar{10}{\downbow}\lsfz{'B}\hup{B}\enotes
\bar
\Notes\lsfz{'B}\hup{B}\enotes
\bar
\Notes\lsfz{'B}\hup{B}\enotes
\bar
\Notes\zchar{10}{\upbow}\lsfz{'B}\hup{B}\enotes
\bar
\Notes\lsfz{'B}\hup{B}\enotes
\bar
\Notes\lsfz{'B}\hup{B}\enotes
\bar
\Notes\zchar{10}{\downbow}\usfz{'E}\hlp{E}\enotes
\bar
\Notes\zchar{10}{\upbow}\usfz{'E}\hlp{E}\enotes
\bar
\Notes\usfz{'E}\hlp{E}\enotes
\bar
\Notes\usfz{'E}\hlp{E}\enotes
\bar
\notes\zchar{-7}{\ff} \zchar{10}{\downbow}\usfz{'E}\cl{E}\ds\enotes
\Notes\qp\qp\enotes
\mulooseness=1
\setdoublebar
\endpiece
\end{music}

%\documentclass{jarticle}
\usepackage{musixdoc}
\startmuflex\makeindex

\begin{document}

\subsubsection*{���㥤���ե�����: �������5�� ��ûĴ ��4�ھϤ��}
\begin{music}
\nostartrule
\setclef1{\bass}
\generalsignature{2}    
\generalmeter{\allabreve}
\parindent 0pt
\startbarno=241
\def\writebarno{\tenrm\the\barno\barnoadd}
\def\raisebarno{2\internote}
\def\shiftbarno{0.1\Interligne}
\systemnumbers
\startpiece\bigaccid
\notes\zchar{18}{\bf Poco pi\`{u} animato}\enotes
\Notes\zchar{-5}{\ff}\ql{'E!e^de}\enotes
\bar
\Notes\ql{ba'GF}\enotes
\bar
\Notes\ql{'E!e^de}\enotes
\bar
\Notes\ql{ba'GF}\enotes
\bar
\Notes\ql{'E!e^de}\enotes
\bar
\Notes\ql{ba'GF}\enotes
\bar
\Notes\ql{'E!e^de}\enotes
\bar
\Notes\ql{ba'GF}\enotes
\bar
\Notes\ql{'E}\qp\enotes
\NOtes\zchar{-5}{\fff}\hu{'B}\enotes
\setdoublebar
\endpiece
\end{music}

\endmuflex
\end{document}

%\subsubsection*{ベートーヴェン: 交響曲第9番 ニ短調 「合唱付き」 第4楽章より}
\begin{music}
\nostartrule
\setclef1{\bass}
\generalsignature{2}    
\generalmeter{\meterC}
\parindent 0pt
\startbarno=241
\def\writebarno{\tenrm\the\barno\barnoadd}
\def\raisebarno{2\internote}
\def\shiftbarno{0.1\Interligne}
\systemnumbers
\startpiece\bigaccid
\notes\zchar{-5}{\p}\zchar{18}{\bf Allegro assai}\enotes
\Notes\ql{'FFG!a}\enotes
\bar
\Notes\ql{!a'GFE}\enotes
\bar
\Notes\ql{'DDEF}\enotes
\bar
\Notes\ql{'F}\qp\ql{E}\qp\enotes
\bar
\Notes\ql{'FFG!a}\enotes
\bar
\Notes\ql{!a'GFE}\enotes
\bar
\Notes\ql{'DDEF}\enotes
\bar
\Notes\ql{'E}\qp\ql{D}\qp\enotes
\bar
\Notes\ql{'EEFD}\enotes
\bar
\Notes\ql{'E}\enotes
\notes\ibl{0}{'F}{-2}\qb{0}{F}\tbl{0}\qb{0}{G}\enotes
\Notes\ql{'FD}\enotes
\bar
\Notes\ql{'E}\enotes
\notes\ibl{0}{'F}{-2}\qb{0}{F}\tbl{0}\qb{0}{G}\enotes
\Notes\ql{'FE}\enotes
\bar
\Notes\ql{'DE}\qu{A}\ql{F}\enotes
\bar
\Notes\ql{'FFG!a}\enotes
\bar
\Notes\ql{!a'GFE}\enotes
\bar
\Notes\zchar{-5}{\p}\ql{'DDEF}\enotes
\bar
\Notes\ql{'E}\qp\ql{D}\qp\enotes
\setdoublebar
\endpiece
\end{music}

%\subsubsection*{グリーグ: 劇音楽「ペール・ギュント」より「山の魔王の宮殿にて」}
\begin{music}
\nostartrule
\setclef1{\bass}
\generalsignature{2}    
\generalmeter{\meterC}
\parindent 0pt
\startbarno=2
\def\writebarno{\tenrm\the\barno\barnoadd}
\def\raisebarno{2\internote}
\def\shiftbarno{0.1\Interligne}
\systemnumbers
\startpiece\bigaccid
\notes\zchar{18}{\bf Alla marcia e molto marcato}\enotes
\notes\zchar{-5}{\p}\ibu{0}{'B}{2}\qb{0}{BCD}\tbu{0}\qb{0}{E}\enotes
\notes\ibl{0}{'F}{-2}\qb{0}{F}\tbl{0}\qb{0}{D}\enotes
\Notes\ql{'F}\enotes
\bar
\notes\ibl{0}{'E}{-2}\qb{0}{^E}\tbl{0}\qb{0}{^C}\enotes
\Notes\ql{'E}\enotes
\notes\ibl{0}{'E}{-2}\qb{0}{=E}\tbl{0}\qb{0}{=C}\enotes
\Notes\ql{'E}\enotes
\bar
\notes\ibu{0}{'B}{2}\qb{0}{B^CD}\tbu{0}\qb{0}{E}\enotes
\notes\ibl{0}{'F}{1}\qb{0}{FDF}\tbl{0}\qb{0}{!b}\enotes
\bar
\notes\ibl{0}{'F}{0}\qb{0}{!a'FD}\tbl{0}\qb{0}{F}\enotes
\Notes\ql{a}\qp\enotes
\bar
\Notes\qu{'B}\ql{F}\qu{B}\ql{F}\enotes
\bar
\Notes\qu{'B}\ql{F}\qu{B}\ql{F}\enotes
\bar
\Notes\qu{'B}\ql{F}\qu{B}\ql{F}\enotes
\bar
\Notes\ql{'D}\ql{!a}\ql{'D}\ql{!a}\enotes
\setdoublebar
\endpiece
\end{music}

